\chapter{格林函数}
\renewcommand{\currentchapterpath}{chapters/ch03/}

\begin{flushright}
\textit{“运筹策帷帐之中，决胜于千里之外。” ——《史记·高祖本纪》
\footnote{格林函数正是“运筹”工具：不直接写出全部波函数，却能决出谱与响应等关键胜负手。}}
\end{flushright}

格林函数（Green's function）是多体理论的核心对象之一。本章把 Mahan 的谱表示/频率求和与 Gross《Many-Particle Theory》的绘景、Gell-Mann--Low、Rayleigh--Schr\"odinger、Goldstone 连锁集群、粒子--空穴、极化传播子、$G_2$、GMB、骨架自洽、准粒子与 Bethe--Salpeter 合在一条链上。


\chapterinput{pictures}
\chapterinput{green}
\chapterinput{species}
\chapterinput{matsubara}
\chapterinput{fdt}
\chapterinput{eom}
\chapterinput{polarization}
\chapterinput{gellmann_low}
\chapterinput{rayleigh}
\chapterinput{particle_hole}
\chapterinput{wick}
\chapterinput{feynman}
\chapterinput{diagram_rules}
\chapterinput{vacuum_amplitude}
\chapterinput{gmb}
\chapterinput{dyson}
\chapterinput{skeleton}
\chapterinput{g_covariant}
\chapterinput{quasiparticle}
\chapterinput{bethe_salpeter}
\chapterinput{finite_T_pt}
\chapterinput{indep_boson}
\chapterinput{elph_sigma}
\chapterinput{noneq}
